\section{Back Propagation}

To actually train and get the wanted behavior from a MLP, 
we must obtain the direction of largest growth or smallest 
otherwise known as the gradient of the parameters. With the gradient,
we can use that to move in a direction of lower loss. With a little bit
of multi-variable calculus and chain rule, we can obtain all of the local derivatives
and in turn, their gradients. Here is the math behind it.

The most rudimentary case is to imagine taking derivatives at the root.
\begin{equation}
    \frac{\partial L}{\partial c} = \cancelto{1}{\frac{\partial L }{\partial L }} \frac{\partial L}{\partial c} = \frac{\partial L}{\partial c}
\end{equation}

However, imagine if we have a function of two variables with one being a composite.
In this example, it should become clear that we will always calculate the outer derivative and multiply it by the inner derivative.
At the start, we will simply take it's derivative and multiply it by one. 

\begin{definition}[Chain Rule]
    With that, back propagation is defined by the chain rule.
    \begin{align}
        & L(c, d) = c + d, \quad  c(a) = a+b; \\
        \Longrightarrow  & \boxed{\frac{\partial L}{\partial a} = \frac{\partial L }{\partial c } \frac{\partial c}{\partial a}}
    \end{align}

    Acculumating gradients:
    \begin{align}
        & L(c, d) = c(a) + d(a), \quad c(a) = a + b, \ d(a) = a + c\\
       \Longrightarrow & \boxed{\frac{\partial L}{\partial a} =\frac{\partial L}{\partial c}\frac{\partial c}{\partial a}+\frac{\partial L}{\partial d}\frac{\partial d}{\partial a}}
    \end{align}
\end{definition}

\begin{definition}[Gradient]
    We can take this farther and define the gradient of the loss function.
        \begin{equation}
        \boxed{\nabla L = \left\langle \frac{\partial L}{\partial a_0}, \dots \frac{\partial L}{\partial a_i} \right\rangle}
    \end{equation}
\end{definition}

This is a very powerful tool and the next thing to tackle is how we order our function 
so that we can create an algorithm to start from the root and propagate back.

\newpage 

\subsection{Topological Sort}

With back propagation, we will need to be able to know what order to propagate. 
This is basically a job scheduling problem so topological sort will be the perfect algorithm for this.
Topological sort works by applying a post-order DFS and reversing the list. This can easily be done 
using a recursive formula to completely search each branch. Reversing should be extremely simple with a library.

\begin{algorithm}
\caption{Topological Sort}\label{alg:cap}
    \begin{algorithmic}
        \Require $n \geq 0$
        \Ensure $y = x^n$
        \State $y \gets 1$
        \State $X \gets x$
        \State $N \gets n$
        \While{$N \neq 0$}
            \If{$N$ is even}
                \State $X \gets X \times X$
                \State $N \gets \frac{N}{2}$  \Comment{This is a comment}
            \ElsIf{$N$ is odd}
                \State $y \gets y \times X$
                \State $N \gets N - 1$
            \EndIf
        \EndWhile
    \end{algorithmic}
\end{algorithm}